\scanpage{3}

\mspage{1}

Catégories de modèles.

1. \deleted{\unsure{2|3}}\uline{Couples de Quillen.} $(\Phi,\Psi)$.

Je veux d'abord donner une variante affaiblie de la notion de triple de Quillen ( ``Model Category'' de Quillen). \unsure{On|Ce} \unsure{bien|ne} \unsure{sûr|sera} \unsure{pas} une notion tout-à-fait naturelle -- aussi il y aura \uline{deux} variantes affaiblies, duales l'une de l'autre.

Mais avant \unsure{ceci|ça|cela}, un \unsure{sorite|suite} sur les LLP (left lifting property) et RLP (right lifting property). Soit $M$ une catégorie quelconque (le plus souvent, $M$ sera grande). On considère des parties

% slanted
\margin{donner définition en forme, ou renvoi à Quillen}

\[ \Phi,\Psi\subset\mathrm{Fl}(M) \]

et on définit des applications

\[ \boxed{\begin{aligned}\Phi&\longmapsto\Phi_*\\ \Psi&\longmapsto\Psi^*\end{aligned}}\qquad \mathrm{Fl}(M)\longrightarrow\mathrm{Fl}(M) \]

par

\[ \begin{aligned}\Phi_*={}&\{\text{ens. des }p\in\mathrm{Fl}(M)\text{ qui ont \unsure{la|le} RLP p.~r.}\\ &\deleted{\unsure{\text{a}|\text{à}}}\text{à }\Phi\}\\
\Psi^*={}&\{\text{ens des }i\in\mathrm{Fl}\,M\text{ qui ont la LLP p.~r}\\ &\text{à }\Psi\},\end{aligned} \]

et ensuite

\[ \boxed{\begin{aligned}\Phi&\longmapsto\widetilde\Phi\\ \Psi&\longmapsto\overline\Psi\end{aligned}}\qquad \mathrm{Fl}(M)\to\mathrm{Fl}(M) \]

par

\[ \begin{aligned}\widetilde\Phi&=(\Phi_*)^*\\ \overline\Psi&=(\Psi^*)_*\end{aligned} \]

% archive sheet number: 2

